\begin{abstract}
프런티어 LLM 계획기는 처음 보는 장면에서도 과제를 이해하지만 한 번 답하는 데 수 초가 걸리고 호출마다 비용이 든다. 과제 데이터로 미세조정한 VLA는 빠르지만 부정확하고, 환경이 바뀌면 크게 무너진다. 우리는 느린 상위 계획기와 빠른 VLA로 된 계층 PaceNotes를 제안하고, 상위부터 단계적으로 검증한다. 이름은 랠리에서 왔다. 코드라이버는 앞 구간의 노트를 읽어 주고, 결과를 확인하며, 다음 노트를 부른다. PaceNotes에서 명령 권한은 모두 코드라이버인 상위 계획기에 있다. 드라이버인 VLA는 좁은 조이스틱으로, 상위가 허용한 범위 안의 보정, 쥐는 시점, 작은 회피, 부딪힌 뒤의 재정렬과 이상 신호만 맡는다. 3--8B VLM 하나에 계획부터 움직임까지 모두 맡기는 대신 그 짐을 덜어 주는 것이 구조의 핵심 가설이다. 1단계에서는 상위가 단순한 결정적 실행기를 직접 지시하는 상위 단독 구성을 먼저 재고, VLA 결합은 그 뒤에 잰다. 주 기여는 학습 가능한 크기의 상위 계획기다. 무료 시뮬 참값 라벨과 공개 로봇 데이터의 점 라벨로 Qwen3.5-35B-A3B를 LoRA 학습한다. 모델은 거리를 m 숫자로 말하지 않고 머리 영상 위의 점과 높이 의도만 내며, 목표 xyz는 코드가 센서 깊이로 계산한다(줄기 D). 근거는 이렇다. Astra에 맞춘 단독 조종은 \tentative{7/7}을 성공했지만 영샷 오픈급 모델은 대상 위치를 읽지 못했다(\tentative{0/5}, 대리 \tentative{0/4}). 시뮬 참값으로 학습한 8B는 접근 목표 오차를 \tentative{125.5\,mm}에서 \tentative{2.6\,mm}로 줄였고, 35B는 지연 p50 \tentative{1.04\,s}로 분포 안 폐루프 \tentative{4/4}였다. 그러나 한 높이 데이터로는 크기와 무관하게 본 적 없는 탁자 높이에서 무너졌고, 깊이 점 출력만 새 높이에서 \tentative{3.5\,mm}를 지켰다. 사람이 적어 준 탁자 높이는 참값이라 실제에 없으므로 넣지 않고, 높이는 깊이가 맡는다. 다양화 시뮬로 학습한 35B 소규모 비교에서 줄기 D는 새 물체 범주에서 \tentative{4.0\,mm}(하이브리드 \tentative{83.7\,mm}), 폐루프 \tentative{6/8}로 채택되었고, 이 개선은 공개 점 데이터 덕이었다(빼면 \tentative{85.3\,mm}). 8B 방향 실험에서 공개/시뮬 비율을 약 0.6보다 높여도 이득이 없었고, 사실적 시뮬 L8S는 많을수록 나았으나 50\% 뒤로는 이득이 작았다. L8S 1,000편 확인판은 새 물체에서 크게 나아졌지만(중앙 \tentative{40.9$\rightarrow$9.6\,mm}) 옛 시뮬 평가와 3인칭 공개 평가에서 나빠졌고, 원인 진단 뒤 사용자 결정으로 본 학습을 시작했다. 본 35B는 L8S 139,603행과 공개 점 62,794행으로 3 에폭 학습 중이며, 0.5 에폭에서 옛 시뮬 평가의 실패율 차가 줄어드는 것이 보인다(\tentative{초기값}). 다음 데이터는 L8S의 5배 규모·다양성의 생성기 L9이고, VLA 결합은 결합 손실 5\%p 이하를 채택 기준으로 하고, 짐을 덜어 준 VLA와 단독 VLA를 직접 비교한다.
\end{abstract}
